\documentclass[../../main.tex]{subfiles}
\begin{document}
\begin{oframed}
  \subfile{problem}
\end{oframed}

{\bf [解答]} $\vec{a}=(1,\sqrt{3})$とおく．まず，$t=0$におけるPの速度を表すベクトルが$\vec{b}=(1,0)$だとして考える．$x$軸正方向の移動量を$x$，$y$軸正方向の移動量を$y$とすると，題意の条件から
\begin{align}
  |x|+|y|\le 2
\end{align}
である．これを図示して下図斜線部(境界含む)．

\begin{zahyou}[ul=10mm](-2.5,2.5)(-2.5,2.5)
  \def\Fx{-1*X+2}
  \def\Gx{X+2}
  \def\Hx{-1*X-2}
  \def\Mx{X-2}
  \YGurafu\Fx{0}{2}
  \YGurafu\Gx{-2}{0}
  \YGurafu\Hx{-2}{0}
  \YGurafu\Mx{0}{2}
  \YNurii*\Fx\Mx{0}{2}
  \YNurii*\Gx\Hx{-2}{0}
  \Put{(2,0)}[ne]{$2$}
  \Put{(0,2)}[ne]{$2$}
  \Put{(-2,0)}[nw]{$-2$}
  \Put{(0,-2)}[se]{$-2$}
\end{zahyou}

これを$\vec{b}$が$\vec{a}$の方向にくるまで回転した領域が求める領域で，図示して下図斜線部(境界含む)．
$\cdots$(答)

\begin{zahyou}[ul=10mm](-2.5,2.5)(-2.5,2.5)
  \tenretu*<perl>{%
  A(1,sqrt(3));B(-sqrt(3),1);C(-1,-sqrt(3));D(sqrt(3),-1)}
  \Drawline{\A\B\C\D\A}
  \Put\A[syaei=xy,xlabel=1,ylabel=\sqrt{3}]{}
  \Nuritubusi*{\A\B\C\D}
\end{zahyou}

\newpage
\end{document}
